% GENERATED FILE. Edit canonical lessons, then run npm run generate:books.
% canonical-source-hash: fnv1a64:0bc201fa265a3cd5
% canonical-lessons: AR-C25-yawm

\chapter{Day}
\label{ch:ar-day}
\hlchaptermodality{25}

\begin{chapteropening}
\textbf{By the end of this chapter:} \emph{I can say 'day' in Arabic and recognise its cognate in Hebrew's Yom Kippur.}

Say yawm, run it across the Semitic family, and mark the Egyptian sun-word origin as unsettled.
\end{chapteropening}

\section[\texorpdfstring{\ar{يوم} (yawm)}{يوم (yawm)}]{\ar{يوم} (yawm) — \textquotedblleft{}day,\textquotedblright{} a word shared across the whole Semitic family}

\label{lesson:AR-C25-yawm}



\begin{quote}
This word is bigger than Arabic — it's one of the oldest, most widely shared words across the entire Semitic language family.
\end{quote}

\subsection*{You'll want to know — \ar{يوم} — \textquotedblleft{}day\textquotedblright{}}
\textbf{\ar{يوم}} (\textbf{yawm}) — \textquotedblleft{}\textbf{day}\textquotedblright{} — root \textbf{y-w-m}. Arabic's day-names can compound with \emph{yawm} — \textbf{\ar{يوم} \ar{الجمعة}} (\emph{yawm al-jumʿa}, \textquotedblleft{}day of the gathering\textquotedblright{}) is the full, formal way to say \textquotedblleft{}Friday,\textquotedblright{} the exact same \textquotedblleft{}day-of-X\textquotedblright{} construct-state pattern as Latin's \emph{diēs Lūnae} — but colloquially, the bare day-name alone (\emph{al-jumʿa}, \emph{as-sabt}, already taught earlier) is just as common. \emph{Yawm} itself is also the ordinary, standalone word for \textquotedblleft{}a day.\textquotedblright{}

\begin{culture}
\textbf{Yawm} traces to \textbf{Proto-Semitic} \emph{\textbf{*yawm-}} — and this isn't just an Arabic word with foreign relatives, it's a genuinely \textbf{Common Semitic} word, found with the same shape and meaning across the whole family: Hebrew \textbf{\he{יוֹם}} (\emph{yōm} — yes, the \emph{yom} in \textbf{Yom Kippur}, \textquotedblleft{}Day of Atonement\textquotedblright{}), Akkadian \textbf{ūmum}, and Aramaic \textbf{yawmā}. Four different Semitic languages, one shared ancient word for \textquotedblleft{}day.\textquotedblright{}
\end{culture}

\begin{culture}
Links beyond the Semitic family have been proposed, but none is established. A suggested tie to an Egyptian word for \textquotedblleft{}sun\textquotedblright{} does not hold up: Egyptian's ordinary word for \textquotedblleft{}day\textquotedblright{}, \emph{ḥrw}, is unrelated. The Semitic cognates above are the solid part of this word's story.
\end{culture}

\subsection*{Your turn}
Say these aloud:

\begin{itemize}
\raggedright
  \item \textquotedblleft{}yawm\textquotedblright{} — \textquotedblleft{}day\textquotedblright{}
  \item the Semitic family — yawm, yōm (Yom Kippur), ūmum, yawmā
  \item the honest limit — the Semitic family is solid; links beyond it are not established
\end{itemize}

\begin{tcolorbox}[breakable,colback=teal!4,colframe=teal!35!black,title={Before you move on}]
What Hebrew word, familiar from a well-known holiday name, is \textbf{yawm}'s direct cognate? (\textbf{Yōm}, as in \textbf{Yom Kippur}.) Name two other Semitic languages that share this exact word for \textquotedblleft{}day.\textquotedblright{} (\textbf{Akkadian} \emph{ūmum}, \textbf{Aramaic} \emph{yawmā}.) Is any link beyond the Semitic family established? (\textbf{No} — the Semitic cognates are the solid part.)
\end{tcolorbox}
